\documentclass[11pt]{article}
\usepackage[french]{babel}
\usepackage[utf8]{inputenc}
\usepackage[T1]{fontenc}
\usepackage{xspace}
\usepackage{array}
\usepackage{color}
\usepackage{listings}
\usepackage{fancyvrb} %verbatim
\usepackage{tabularx}
\usepackage{eurosym}
\usepackage{mathrsfs}
\usepackage{stmaryrd}
\usepackage{amssymb}
\usepackage[?]{amsmath}
\usepackage{epsfig}
\usepackage{float}
\usepackage{wrapfig}
\usepackage{pifont}
\usepackage{enumerate}
\usepackage[np]{numprint} %\np{12345} donne 12 345 
\usepackage{multicol}
\newcommand{\second}{2\up{d}\xspace}
\newcommand{\seconde}{2\up{de}\xspace}
\newcommand{\R}{\mathbb R}
\newcommand{\Rp}{\R_+}
\newcommand{\Rpe}{\R_+^*}
\newcommand{\Rm}{\R_-}
\newcommand{\Rme}{\R_-^*}
\newcommand{\N}{\mathbb N}
\newcommand{\D}{\mathbb D}
\newcommand{\Q}{\mathbb Q}
\newcommand{\Z}{\mathbb Z}
\newcommand{\C}{\mathbb C}
\newcommand{\grs}{\mathfrak S}
\newcommand{\IN}[1]{\llbracket 1,#1\rrbracket}
\newcommand{\card}{\text{Card}\,}
\usepackage{mathrsfs}
\newcommand{\parties}{\mathscr P}
\renewcommand{\epsilon}{\varepsilon}
\newcommand{\rmd}{\text{d}}
\newcommand{\diff}{\mathrm D}
\newcommand{\Id}{{\rm Id}}
\newcommand{\e}{{\rm e}}
\newcommand{\I}{{\rm i}}
\newcommand{\J}{{\rm j}}
\newcommand{\ro}{\circ}
\newcommand{\exu}{\exists\,!\,}
\newcommand{\telq}{\,\, \mid \,\,}
\newcommand{\para}{\raisebox{0.1em}{\text{\footnotesize /\hspace{-0.1em}/}}}   
\newcommand{\vect}[1]{\overrightarrow{#1}}
\newcommand{\scal}[2]{\left(\, #1 \mid #2 \, \right)}
\newcommand{\ortho}[1]{{#1}^\perp}
\newcommand{\veci}{\vec{\text{\it \i}}}
\newcommand{\vecj}{\vec{\text{\it \j}}}
\newcommand{\rep}{$(O;\veci,\vecj,\vec{k})$\xspace}
\newcommand{\Oijk}{$(O, \veci,\vecj,\vec{k})$\xspace}
\newcommand{\rond}{repère orthonormal direct}
\newcommand{\bond}{base orthonormale directe}
\newcommand{\eq}{\Longleftrightarrow}
\newcommand{\implique}{\Longrightarrow}
\newcommand{\noneq}{\ \ \ /\hspace{-1.45em}\eq}
\newcommand{\tend}{\longrightarrow}
\newcommand{\egx}[2]{\underset{#1 \tend #2}=}
\newcommand{\asso}{\longmapsto}
\newcommand{\vers}{\longrightarrow}
\newcommand{\eqn}{~\underset{n \rightarrow \infty}{\sim}~}
\newcommand{\eqx}[2]{~\underset{#1 \rightarrow #2}{\sim}~}
\newcommand{\egn}{~\underset{n \rightarrow \infty}{=}~}
\renewcommand{\descriptionlabel}{\hspace{\labelsep}$\bullet$}
\renewcommand{\bar}{\overline}
\DeclareMathOperator{\ash}{{Argsh}}
\DeclareMathOperator{\cotan}{{cotan}}
\DeclareMathOperator{\ach}{{Argch}}
\DeclareMathOperator{\ath}{{Argth}}
\DeclareMathOperator{\sh}{{sh}}
\DeclareMathOperator{\ch}{{ch}}
%\DeclareMathOperator{\tanh}{{th}}
\DeclareMathOperator{\Mat}{{Mat}}
\DeclareMathOperator{\Vect}{{Vect}}
\DeclareMathOperator{\trace}{{tr}}
\newcommand{\tr}{{}^{\mathrm t}}
\newcommand{\divi}{~\big|~}
\newcommand{\ndivi}{~\not{\big|}~}
\newcommand{\et}{\wedge}
\newcommand{\ou}{\vee}
%\DeclareMathOperator{\trace}{{trace}}
\renewcommand{\det}{\operatorname{\text{dét}}}
%\DeclareMathOperator{\det}{{dét}}
\DeclareMathOperator{\grad}{{grad}}
%\DeclareMathOperator{\th}{th}
%\renewcommand{\tanh}{\mathop{\mathrm{th}}{}}
\renewcommand{\arcsin}{{\mathop{\mathrm{Arcsin}}}}
\renewcommand{\arccos}{{\mathop{\mathrm{Arccos}}}}
\renewcommand{\arctan}{{\mathop{\mathrm{Arctan}}}}
%\DeclareMathOperator{\arctan}{{Arctan}}
\renewcommand{\tanh}{{\mathop{\mathrm{th}}}}
\newcommand{\pgcd}{\mathop{\mathrm{pgcd}}}
\newcommand{\ppcm}{\mathop{\mathrm{ppcm}}}
\newcommand{\fonc}[4]{\left\{\begin{tabular}{ccc}$#1$ & $\vers$ & $#2$ \\
$#3$ & $\asso$ & $#4$ \end{tabular}\right.}
\renewcommand{\geq}{\geqslant}
\renewcommand{\leq}{\leqslant}
\renewcommand{\Re}{\text{\rm Re}}
\renewcommand{\Im}{\text{\rm Im}}
\renewcommand{\ker}{\mathop{\mathrm{Ker}}}
\newcommand{\Lin}{\mathcal L}
\newcommand{\GO}{\mathcal O}
\newcommand{\GSO}{\mathcal{SO}}
\newcommand{\GL}{\mathcal{GL}}
\renewcommand{\emptyset}{\varnothing}
\newcommand{\arc}[1]{\overset{\frown}{#1}}
\newcommand{\rg}{\mathop{\mathrm{rg}}}
\newcommand{\ds}{\displaystyle}
\newcommand{\co}[3]{\begin{pmatrix}#1 \\ #2 \\ #3\end{pmatrix}}
\newcommand{\demi}{\frac 1 2}

\newcommand{\exod}[1]{\addtocounter{cexo}{1} \vskip 1em 
\noindent{\bf \large Exercice \thecexo.\ #1}}
\newcounter{cact}
\renewcommand{\thecact}{\arabic{cact}}
\newcommand{\act}{\addtocounter{cact}{1} \vskip 1em \noindent{\bf Activité \thecact.\ }}
%\renewcommand{\labelitemi}{\m@th$\star$}
\newcommand{\limi}[2]{\underset{#1 \rightarrow #2}\lim}    
\newcommand{\nterme}[1]{{\em #1}}
\newcommand{\dif}[1]{#1}
\renewcommand{\FrenchLabelItem}{\textbullet} %bullet dans les items
%\renewcommand{\labelitemi}{\textbullet}
%\renewcommand{\labelitemii}{\textbullet}
%\renewcommand{\labelitemiii}{\textbullet}
%\renewcommand{\theenumi}{\alph{enumi}}
%\setenumerate[0]{label=\alph}
\usepackage{xr}
\usepackage{hyperref}
\usepackage{textcomp} %pour les ' dans lstlisting
\usepackage{listings}
%------------------ info
\lstset{
  language=Python,
  commentstyle=\color{red},
  upquote=true,
  literate={é}{{\'e}}1
           {ê}{{\^e}}1
           {Ã }{{\`a}}1
           {è}{{\`e}}1
           {î}{{\^i}}1
           {ï}{{\"i}}1
           {ë}{{\"e}}1
           {ù}{{\`u}}1,  
  columns=flexible,
  basicstyle=\ttfamily,
  keywordstyle=\color{blue},
  %identifierstyle=\color{green},
  stringstyle=\ttfamily,
  showstringspaces=false,
  %numbers=left, 
  %numberstyle=\tiny, 
  %stepnumber=1, 
  %numbersep=5pt,
  frame=L
}

 
\ifdefined\HCode
  \def\pgfsysdriver{pgfsys-dvisvgm4ht.def}
\fi 
\usepackage{pgf,tikz,pgfplots}
\usepackage{tkz-tab,tkz-fct}
\pgfplotsset{compat=1.16}
%\usepackage{mathrsfs}
\usetikzlibrary{arrows}

%-----------------------
\usepackage{fancybox}
\usepackage{fancyhdr}

\definecolor{gray}{HTML}{c8caca}
\definecolor{bleusad}{HTML}{0056B3}
\newcommand{\titreFiche}[3]{\centerline{\Ovalbox{\parbox[t]{0.7\textwidth}{
\vspace{-0.5em}\begin{center} \bf \large #1
\end{center}
\vspace{-0.5em}}}}

  \bigskip
\fancyfoot[L]{\includegraphics[height=1em]{/var/www/sacado/static/img/sacadoA1.png}\textcolor{bleusad}{\textsf{\ \raisebox{2pt}{ {{\footnotesize https://sacado.xyz > #3}}}}}}
\fancypagestyle{Page1}{\fancyhead{}\renewcommand{\headrulewidth}{0pt}}
\thispagestyle{Page1}
\fancyhead[L]{#1}
\fancyfoot[R]{\textsc{#2}}
\renewcommand{\footrulewidth}{0.4pt}
\pagestyle{fancy}
}

\newcounter{cexo}
\renewcommand{\thecexo}{\arabic{cexo}}
\newcommand{\exo}{\refstepcounter{cexo} \vskip 1.5em \noindent{\bf Exercice \thecexo.\ }}



\definecolor{bleu2}{rgb}{0,0.58,0.79}
\definecolor{violet}{rgb}{0.79,0.66,0.98}
\definecolor{gray}{rgb}{0.79,0.66,0.98}

\newcommand{\exercice}[1]{ % Exercice directement dans la page
\vspace{0.2cm}
\tikz\node[rounded corners=2pt, fill=bleu2]{\color{white}\textbf{#1}};
\vspace{0.1cm}
}

\newcommand{\sacado}[1]{ % Exercice directement dans la page
\vspace{0.2cm}
\tikz\node[rounded corners=2pt, fill=violet]{\color{white}\textbf{#1}};
}


\newcommand{\competence}[1]{
  \hfill\textcolor{bleusad}{#1}
}
\newcommand{\competences}[1]{
  {\bf Compétences ciblées : }
  \begin{enumerate}
   \item   #1
  \end{enumerate}

  \medskip
}

\newcommand{\savoir}[1]{
  {\bf Autre savoir-faire ciblé : } 
  #1
  \medskip
}

\newcommand{\savoirs}[1]{
  %  {\bf Autres savoir-faire ciblés : }
 \vspace{-1em}
  \textcolor{bleusad}{\it \begin{itemize}
      #1
  \end{itemize}
  \centerline{\rule[0.5em]{5cm}{0.5pt}}}
  \medskip
}


\newcommand{\entreexo}{ % entre les exos}
  \medskip
  \centerline{\hrule{5cm}{1pt}}
  \medskip
}

\setlength{\parindent}{0cm}
\usepackage[left=1cm,right=1cm,top=0.1cm,bottom=1.5cm]{geometry}

\usepackage{multicol}

%%%%% Pour faire des boites
\usepackage[tikz]{bclogo}
\usepackage{bclogo}
\usepackage{framed}
\usepackage[skins]{tcolorbox}
\tcbuselibrary{breakable}
\tcbuselibrary{skins}
\usetikzlibrary{quotes,babel,arrows.meta,shadows,decorations.pathmorphing,decorations.markings,patterns}
\usepackage{tikzpagenodes}
\usetikzlibrary{plotmarks}


\usepackage{multido}
\newcommand{\point}[1]{\vspace{0.1cm}\multido{}{#1}{ \dotfill \medskip \endgraf}}
\newcommand{\ligne}[1]{\vspace{0.1cm}\multido{}{#1}{ {\color{cqcqcq}\hrulefill} \medskip \endgraf}}

\newcounter{exo}
\newcounter{cptr}




\renewenvironment{leftbar}[1][\hsize]
{% 
    \def\FrameCommand
    {%
       {\color{sacado_purple}\vrule width 1pt}%
        \hspace{0pt}%must no space.
        \fboxsep=\FrameSep%\colorbox{yellow}%
    }%
    \MakeFramed{\hsize#1\advance\hsize-\width\FrameRestore}%
  }
{\endMakeFramed}


\newenvironment{leftbarPink}[1][\hsize]
{% 
    \def\FrameCommand
    {%
       {\color{sacado_pink}\vrule width 1pt}%
        \hspace{0pt}%must no space.
        \fboxsep=\FrameSep%\colorbox{yellow}%
    }%
    \MakeFramed{\hsize#1\advance\hsize-\width\FrameRestore}%
  }
{\endMakeFramed}


\newcommand{\headerGeneral}[1]{ % intitulé, couleur, qrcode
\begin{tikzpicture}[remember picture,overlay]
\coordinate(NO) at (-1.5,0);
\coordinate(SW) at (22,1);
\node[text width=6.2cm,font=\small] at (20.7,-0.04) { \color{white}{#1}  }; 
\end{tikzpicture}
}


%%%%%% Environnement "pageCours" %%%%%%
\newenvironment{pageIntro}{
\rhead{}
\lhead{}%
}
{\newpage} 

%%%%%% Environnement "pageCours" %%%%%%
\newenvironment{pageCours}{
\rhead{\includegraphics[scale=1]{/var/www/sacado/ressources/tex/pageCours.jpg} }
\lhead{}%
\headerGeneral{p/1234}
}
{\newpage}

%%%%%% Environnement "pageAD" %%%%%%
\newenvironment{pageAd}{
\rhead{\includegraphics[scale=1]{/var/www/sacado/ressources/tex/pageAd.jpg} }
\lhead{}%
\headerGeneral{p/1234}
}
{\newpage}

%%%%%% Environnement "pageParcoursu" %%%%%%
\newenvironment{pageParcoursu}{

\rhead{\includegraphics[scale=1]{/var/www/sacado/ressources/tex/pageParcoursu.jpg} }
\lhead{}%
\headerGeneral{p/1234}
}
{\newpage}
%%%%%% Environnement "pageParcoursd" %%%%%%
\newenvironment{pageParcoursd}{
\rhead{\includegraphics[scale=1]{/var/www/sacado/ressources/tex/pageParcoursd.jpg} }
\lhead{}%
\headerGeneral{p/1234}
}
{\newpage}

%%%%%% Environnement "pageParcourst" %%%%%%
\newenvironment{pageParcourst}{
\rhead{\includegraphics[scale=1]{/var/www/sacado/ressources/tex/pageParcourst.jpg} }
\lhead{}%
\headerGeneral{p/1234}
}
{\newpage}
 
 

%%%%%% Environnement "pageAuto" %%%%%%
\newenvironment{pageAuto}{
\rhead{\includegraphics[scale=1]{/var/www/sacado/ressources/tex/pageAuto.jpg} }
\lhead{}%
\headerGeneral{p/1234}
}
{\newpage}

%%%%%% Environnement "pageAlgo" %%%%%%
\newenvironment{pageAlgo}{
\rhead{\includegraphics[scale=1]{/var/www/sacado/ressources/tex/pageAlgo.jpg} }
\lhead{}%
\headerGeneral{p/1234}
}
{\newpage}
 
 

%%%%%% Numéro de page %%%%%%
\fancyfoot[L]{\colorbox{sacado_purple!70}{\color{white}\thepage}}
\fancyfoot[C]{}

%%%%%% Sections et numéros de sections %%%%%%
 
\renewcommand{\thesection}{\arabic{section}}



\titleformat{\section}{}{%
\hspace{-0.15cm}{%
 \tikz[baseline=(a.base)]\node[draw,circle,inner sep=2pt, line width=0mm, outer sep=2pt,fill=sacado_purple!70](a){\Large {\color{white}\thesection}  };
 }}{1em}{\bf \Large {\color{sacado_purple}#1}}
  
  
\renewcommand{\thesubsection}{\arabic{subsection}}

\titleformat{\subsection}{}{%
\hspace{-0.15cm}{%
 \tikz[baseline=(a.base)]\node[draw,circle,inner sep=2pt, outer sep=2pt,fill=sacado_purple!70](a){\large {\color{white}\thesubsection}  };
 }}{1em}{\bf \large {\color{sacado_purple}#1}}




\makeatletter
\newenvironment{TraitV}[1]{%
% #1 couleur du trait (par défaut CouleurA)
% #2 largeur du trait
% #3 distance entre le trait et le texte
\def\FrameCommand{{\color{#1}\vrule width 2pt}
\hspace{1em}}\MakeFramed {\advance\hsize-\width}}%
{\endMakeFramed}
\makeatother

%----------------------------------------
%
%   Définitions des environnements de Définitions, propriétés...
%
%----------------------------------------
\newenvironment{Def}{%
\begin{tcolorbox}[colback=sacado_blue_light,colframe=sacado_blue_light,fonttitle=\bfseries,
 boxsep=5pt,left=5pt,right=5pt,top=2pt,bottom=2pt,
colbacktitle=sacado_blue,  enhanced,outer arc=0pt, attach boxed title to top left={xshift=-2mm, yshift=-3mm},
title= \textbf{$\searrow$}\; Définition. ]\vspace{3mm}
}%
{%
\end{tcolorbox}\par} 



\newenvironment{Th}{%
\begin{tcolorbox}[colback=sacado_yellow_light,colframe=sacado_yellow_light,fonttitle=\bfseries,
 boxsep=5pt,left=5pt,right=5pt,top=2pt,bottom=2pt,
colbacktitle=sacado_pink,  enhanced,outer arc=0pt, attach boxed title to top left={xshift=-2mm, yshift=-3mm},
title= \textbf{$\searrow$}\; Théorème. ]\vspace{3mm}
}%
{%
\end{tcolorbox}\par} 


\newenvironment{Pp}{%
\begin{tcolorbox}[colback=sacado_yellow_light,colframe=sacado_yellow_light,fonttitle=\bfseries,
 boxsep=5pt,left=5pt,right=5pt,top=2pt,bottom=2pt,
colbacktitle=sacado_pink,  enhanced,outer arc=0pt, attach boxed title to top left={xshift=-2mm, yshift=-3mm},
title= \textbf{$\searrow$}\; Propriété. ]\vspace{3mm}
}%
{%
\end{tcolorbox}\par} 


\newenvironment{Att}{%
\begin{tcolorbox}[colback=sacado_pink_light,colframe=sacado_pink_light,fonttitle=\bfseries,
 boxsep=5pt,left=5pt,right=5pt,top=2pt,bottom=2pt,
colbacktitle=sacado_pink,  enhanced,outer arc=0pt, attach boxed title to top left={xshift=-2mm, yshift=-3mm},
title= \textbf{$\searrow$}\; Attention. ]\vspace{3mm}\begin{leftbarPink}}%
{\end{leftbarPink}%
\end{tcolorbox}\par}



\newenvironment{Rq}{%
\begin{tcolorbox}[colback=sacado_yellow_light,colframe=sacado_yellow_light,fonttitle=\bfseries,
 boxsep=5pt,left=5pt,right=5pt,top=2pt,bottom=2pt,
colbacktitle=sacado_pink,  enhanced,outer arc=0pt, attach boxed title to top left={xshift=-2mm, yshift=-3mm},
title= \textbf{$\searrow$}\; Propriété.  ]\vspace{3mm}
}%
{%
\end{tcolorbox}\par} 



\newenvironment{Rg}{%
\begin{tcolorbox}[colback=sacado_yellow_light,colframe=sacado_yellow_light,fonttitle=\bfseries,
 boxsep=5pt,left=5pt,right=5pt,top=2pt,bottom=2pt,
colbacktitle=sacado_pink,  enhanced,outer arc=0pt, attach boxed title to top left={xshift=-2mm, yshift=-3mm},
title= \textbf{$\searrow$}\; Règle. ]\vspace{3mm}
}%
{%
\end{tcolorbox}\par} %%%%%%%%%%%%% Définitions
%%%%%%%%%%%%%  

 
\newenvironment{lBar} {%
\begin{leftbar}
}%
{\end{leftbar}%
}

\newenvironment{Bloc}{%
\hspace{1mm}
\begin{minipage}{0.99\linewidth}
}%
{\end{minipage}%
}





\newenvironment{lBarPink} {%
\begin{leftbarPink}
}%
{\end{leftbarPink}%
}



\newenvironment{Ex}{%
\hspace{1mm} 
\begin{minipage}{0.97\linewidth} \vspace{1mm}
}%
{\vspace{1mm} \end{minipage} %
}



\newtcolorbox{Mt}{%
enhanced,
breakable,
sharp corners,arc=0pt,outer arc=0pt,
boxrule=2pt,boxsep=0pt,top=4pt,left=4pt,right=4pt,bottom=4pt,middle=0pt,
colback=white, 
colframe=white, 
pad at break=0pt,bottomrule at break=0pt,toprule at break=0pt,
borderline east={1pt}{-0.25pt}{sacado_green,dotted},
borderline west={1pt}{-0.25pt}{sacado_green,dotted},
borderline south={1pt}{-0.25pt}{sacado_green,dotted},
borderline north={1pt}{-0.25pt}{sacado_green,dotted},
fontupper={\color{sacado_green}}
        \fcolorbox{sacado_green}{sacado_green}{\color{white}{\small MÉTHODE}}
        \color{sacado_green}\vspace{1mm}
}


\newtcolorbox{ExR}{%
enhanced,
breakable,
sharp corners,arc=0pt,outer arc=0pt,
boxrule=2pt,boxsep=0pt,top=4pt,left=4pt,right=4pt,bottom=4pt,middle=0pt,
colback=white, 
colframe=white, 
pad at break=0pt,bottomrule at break=0pt,toprule at break=0pt,
borderline east={1pt}{-0.25pt}{sacado_green,dotted},
borderline west={1pt}{-0.25pt}{sacado_green,dotted},
borderline south={1pt}{-0.25pt}{sacado_green,dotted},
borderline north={1pt}{-0.25pt}{sacado_green,dotted},
fontupper={\color{sacado_green}}
        \fcolorbox{sacado_green}{sacado_green}{\color{white}{\small EXERCICE RÉSOLU}}
        \color{sacado_green}\vspace{1mm}
}



%%%%%%%%%%%%% ExoCad 7 paramètres : Compétences, qrcode , calculatrice, python, scratch, tableur, annales
 
\newenvironment{Exo}{% code avant
 
\stepcounter{exo}

 
 }

%%%%%%%%%%%%%%%%%%%%%%%%%%%%%%%%%%%%%%%%%%%%%%%%%%%%%%%%%%%%%%%%%%%%%%%%%%%%%%%%%%%%%%%%%%%%%%%%%%%%%%%%%%%%%%%%%%%%%%
%%%%%%%%%%%%%%%%%%%%%%%%%%%%%%%%%%%%%%%%%%%%%%%%%%%%%%%%%%%%%%%%%%%%%%%%%%%%%%%%%%%%%%%%%%%%%%%%%%%%%%%%%%%%%%%%%%%%%%
%%%%%%%%%%%%%%%%  Exercices   sans contours                            %%%%%%%%%%%%%%%%%%%%%%%%%%%%%%%%%%%%%%%%%%%%%%%
%%%%%%%%%%%%%%%%%%%%%%%%%%%%%%%%%%%%%%%%%%%%%%%%%%%%%%%%%%%%%%%%%%%%%%%%%%%%%%%%%%%%%%%%%%%%%%%%%%%%%%%%%%%%%%%%%%%%%%
%%%%%%%%%%%%%%%%%%%%%%%%%%%%%%%%%%%%%%%%%%%%%%%%%%%%%%%%%%%%%%%%%%%%%%%%%%%%%%%%%%%%%%%%%%%%%%%%%%%%%%%%%%%%%%%%%%%%%%


\newcommand{\Sf}[1]{ \vspace{0.1cm}
{\color{sacado_purple}{\Large \textbf{#1}}  } 
} 



\newcommand{\Sfe}[1]{ \vspace{0.1cm}
{\color{sacado_blue}{\Large \textbf{#1}}  } 
} 
 


\usepackage{makeidx}
\makeindex


 